% KIT & FZI research poster, built on Gemini.
% The values, plots, and prose below are illustrative template content.
% Options: kit-green, kit-blue, kit-black70, fzi-green.
\providecommand{\kitfzipalette}{kit-green}
% Use rgb for the documented screen colors, or cmyk for KIT print values.
\providecommand{\kitfzicolormodel}{rgb}
\documentclass[final]{beamer}
\usepackage[T1]{fontenc}
\usepackage{lmodern}
\usepackage[size=custom,width=120,height=72,scale=1.0]{beamerposter}
\usetheme{kitfzi}
\usepackage{graphicx}
\usepackage{booktabs}
\usepackage{tikz}
\usepackage{pgfplots}
\usepackage{anyfontsize}
\usetikzlibrary{arrows.meta,positioning,calc}
\pgfplotsset{compat=1.18}
% Self-contained, illustrative figures. Replace with your research data.
\pgfplotsset{kitfzi axis/.style={
  axis lines=left,axis line style={kitfzibrandgray},tick style={kitfzibrandgray},
  grid=major,grid style={gray!15},
  label style={font=\small},tick label style={font=\small},
  legend style={font=\small,draw=none,fill=none},
  every axis plot/.append style={line width=2.4pt},clip=true
}}

\newcommand{\ablationplot}{%
  \begin{tikzpicture}
    \begin{axis}[kitfzi axis,width=0.91\linewidth,height=8.3cm,
      ybar,bar width=22pt,ymin=60,ymax=86,
      symbolic x coords={Base,+ A,+ B,Full},xtick=data,ylabel={Score},
      enlarge x limits=0.18,nodes near coords,
      every node near coord/.append style={font=\small,text=black}]
      \addplot[fill=kitfziprimary,draw=kitfziprimary] coordinates {(Base,72.4) (+ A,76.1) (+ B,78.5) (Full,81.3)};
    \end{axis}
  \end{tikzpicture}%
}

\newcommand{\methoddiagram}{%
  \begin{tikzpicture}[
    stage/.style={draw=kitfziprimary,line width=1.5pt,rounded corners=5pt,
      fill=kitfzilight,text width=6.7cm,minimum height=3cm,align=center,font=\small},
    flow/.style={-{Latex[length=6mm]},line width=1.7pt}]
    \node[stage] (input) {\textbf{Input}\\$x$};
    \node[stage,right=1.6cm of input] (representation) {\textbf{Representation}\\$h=f_\theta(x)$};
    \node[stage,right=1.6cm of representation] (output) {\textbf{Prediction}\\$\hat y=g_\phi(h)$};
    \node[stage,below=2cm of representation,fill=white] (component) {\textbf{Proposed component}\\Your main idea};
    \draw[flow] (input) -- (representation);
    \draw[flow] (representation) -- (output);
    \draw[flow,kitfziprimary] (component) -- (representation);
    \draw[flow,kitfzisecondary,dashed] (input.north) -- ++(0,1.6cm) -| (output.north);
    \node[above=1.6cm of representation,font=\small,text=kitfzisecondary] {Baseline comparison};
  \end{tikzpicture}%
}

\newcommand{\learningplot}[3]{%
  \begin{tikzpicture}
    \begin{axis}[kitfzi axis,width=0.84\linewidth,height=8cm,
      xmin=0,xmax=50,ymin=0.35,ymax=1,xlabel={Epoch},ylabel={Score},
      title={#3},title style={font=\small},legend pos=south east,
      domain=0:50,samples=51]
      \addplot[kitfziprimary] {#2-0.48*exp(-x/(10+#1*12))};
      \addplot[kitfzisecondary,dashed] {#2-0.10-0.4*exp(-x/(13+#1*12))};
      \legend{Method,Baseline}
    \end{axis}
  \end{tikzpicture}%
}

\newcommand{\sensitivityplot}{%
  \begin{tikzpicture}
    \begin{axis}[kitfzi axis,width=0.91\linewidth,height=8cm,
      domain=-2:2,xmin=-2,xmax=2,ymin=-1.2,ymax=1.2,
      xlabel={Parameter $\tau$},ylabel={Response},legend pos=north west,samples=61]
      \addplot[kitfziprimary] {0.48*x+0.05*sin(deg(2*x))};
      \addplot[kitfzisecondary] {0.35*x+0.16*sin(deg(2*x))};
      \addplot[kitfzibrandgray,dashed] {0.5*x};
      \legend{Method,Baseline,Reference}
    \end{axis}
  \end{tikzpicture}%
}

\newcommand{\trajectorypanel}[3]{%
  \begin{tikzpicture}
    \begin{axis}[width=\linewidth,height=7.8cm,xmin=0,xmax=5,
      ymin=-0.5,ymax=2.5,hide axis,axis equal image,clip=true,
      title={#3},title style={font=\small},axis background/.style={fill=kitfzicanvas}]
      \addplot[gray!30,line width=19pt] coordinates {(0,0) (2.3,0) (5,2)};
      \addplot[white,dashed,line width=1.3pt] coordinates {(0,0) (2.3,0) (5,2)};
      \addplot[#2,line width=3pt,domain=0.3:4.6,samples=25] {#1*(x/4.6)^2};
      \addplot[only marks,mark=square*,mark size=7pt,#2] coordinates {(0.5,0)};
    \end{axis}
  \end{tikzpicture}%
}


% For N columns: (N+1)*sepwidth + N*colwidth = paperwidth.
\newlength{\sepwidth}
\newlength{\colwidth}
\setlength{\sepwidth}{0.025\paperwidth}
\setlength{\colwidth}{0.3\paperwidth}
\newcommand{\separatorcolumn}{\begin{column}{\sepwidth}\end{column}}

\title{Your research title: a clear statement of the main idea}
\author{First Author\inst{1,2} \and Second Author\inst{1} \and Third Author\inst{2}}
\institute[shortinst]{\inst{1} Karlsruhe Institute of Technology (KIT)\\
  \inst{2} FZI Research Center for Information Technology}
\logoleft{\includegraphics[height=4.4cm]{assets/logos/kit-white.pdf}}
\logoright{\includegraphics[height=5.4cm]{assets/logos/fzi-white.pdf}}
\footercontent{Conference / Workshop \textbullet\ City \textbullet\ Year\hfill
  \href{https://example.com/research}{example.com/research}\hfill
  \href{mailto:name@institution.example}{name@institution.example}}

\begin{document}
\begin{frame}[t]
\vspace{1.2ex}
\begin{columns}[t]
\separatorcolumn

\begin{column}{\colwidth}
  \begin{block}{Research questions}
    \begin{enumerate}
      \item What problem matters, and who benefits from solving it?
      \item Which limitation of existing methods motivates your work?
      \item What evidence would demonstrate a meaningful improvement?
    \end{enumerate}
    Introduce the task and the setting in two or three sentences. State the
    assumptions that define the scope of your research.
  \end{block}

  \begin{block}{Contributions}
    \begin{enumerate}
      \item \textbf{Idea.} Describe the main conceptual contribution in one sentence.
      \item \textbf{Reproducible method.} Explain what others can implement or reuse.
      \item \textbf{Evidence and limitations.} Summarize what the experiments establish.
    \end{enumerate}
  \end{block}

  \begin{block}{Method}
    \textbf{1. Input.} Specify the observations, labels, or signals available to the method.

    \textbf{2. Representation.} Introduce the core transformation and the proposed component.
    \[
      \hat y = g_\phi(h), \qquad h = f_\theta(x)
    \]
    Here, $x$ denotes an input, $h$ its representation, and $\hat y$ a prediction;
    $\theta$ and $\phi$ are learned parameters.

    \textbf{3. Evaluation.} Compare against appropriate baselines using the same protocol.
  \end{block}

  \begin{block}{Ablation study}
    \begin{center}\ablationplot\end{center}
    {\small Illustrative values: replace the scores with your own measurements.}
  \end{block}

  \begin{exampleblock}{Experimental setup}
    \textbf{Data.} State the dataset, split sizes, preprocessing, and access conditions.

    \textbf{Protocol.} Define the metrics, baseline settings, and evaluation procedure.

    \textbf{Reproducibility.} Provide code, configurations, random seeds, and hardware details.
  \end{exampleblock}

  \begin{block}{References}
    {\small Add citations where they support a specific statement; this is a
    formatting example \cite{shannon1948communication}.
    \bibliographystyle{plain}\bibliography{poster}}
  \end{block}
\end{column}

\separatorcolumn
\begin{column}{\colwidth}
  \begin{block}{Overview}
    Use a compact diagram to show the complete method and highlight the part
    introduced by your work.
    \begin{center}\methoddiagram\end{center}
    Read from left to right: input, representation, and prediction. The lower
    component marks the proposed idea; the dashed path identifies the baseline comparison.
  \end{block}

  \begin{block}{Learning curves}
    \begin{center}
      \begin{minipage}{0.48\linewidth}\learningplot{0.2}{0.96}{(a) Training}\end{minipage}\hfill
      \begin{minipage}{0.48\linewidth}\learningplot{0.27}{0.90}{(b) Validation}\end{minipage}
    \end{center}
    {\small Illustrative curves. Solid lines denote the proposed method; dashed lines denote the baseline.}

    Explain how performance changes during training. For real experiments,
    report variation across runs and describe how model selection was performed.
  \end{block}

  \begin{block}{Sensitivity analysis}
    \begin{center}\sensitivityplot\end{center}
    {\small Illustrative response to parameter $\tau$; the dashed gray line is a reference.}

    Identify the range where the method behaves reliably. Explain why the
    parameter matters and where the observed behavior breaks down.
  \end{block}

\end{column}

\separatorcolumn
\begin{column}{\colwidth}
  \begin{exampleblock}{Quantitative results}
    \begin{center}
      \renewcommand{\arraystretch}{1.25}
      \begin{tabular}{lrrr}
        \toprule
        Method & Score $\uparrow$ & Error $\downarrow$ & Time (ms) $\downarrow$ \\
        \midrule
        Baseline & 72.4 & 0.34 & 12.8 \\
        + Component A & 76.1 & 0.29 & 13.1 \\
        + Component B & 78.5 & 0.26 & 13.4 \\
        \textbf{Full method} & \textbf{81.3} & \textbf{0.22} & 13.6 \\
        \bottomrule
      \end{tabular}
    \end{center}
    {\small All values are illustrative. Arrows indicate the preferred metric direction.}

    State the main comparison and its practical significance. Discuss the tradeoff
    between predictive quality and computational cost.
  \end{exampleblock}

  \begin{block}{Qualitative results}
    \begin{center}
      \begin{minipage}{0.31\linewidth}\trajectorypanel{0.05}{kitfzisecondary}{(a) Baseline}\end{minipage}\hfill
      \begin{minipage}{0.31\linewidth}\trajectorypanel{0.65}{kitfziprimary}{(b) Setting A}\end{minipage}\hfill
      \begin{minipage}{0.31\linewidth}\trajectorypanel{1.15}{kitfziprimary}{(c) Setting B}\end{minipage}
    \end{center}
    {\small Schematic trajectory examples for layout illustration; these are not experimental results.}

    Use representative cases to explain what changes. Include a failure case
    when it helps readers understand the method's limitations.
  \end{block}

  \begin{block}{Limitations and next steps}
    \begin{itemize}
      \item Which assumptions constrain the interpretation of the results?
      \item Does the method generalize to different data, settings, or populations?
      \item What next experiment would test the most important open question?
    \end{itemize}
  \end{block}

  \begin{alertblock}{Take-home message}
    \textbf{Write the one sentence you want readers to remember.}

    Connect the main finding to its significance, and keep the conclusion within
    the scope supported by your evidence.
  \end{alertblock}

  \begin{block}{Code and data}
    Help readers reproduce the work: link the code, data, and experiment
    configuration, and state their availability or access conditions.

    \textbf{Project page:} \href{https://example.com/research}{example.com/research}
  \end{block}

  \begin{block}{Acknowledgments}
    Replace this text with the appropriate collaborators, funding acknowledgments,
    and project identifiers. Remove the section if it is not needed.
  \end{block}
\end{column}

\separatorcolumn
\end{columns}
\end{frame}
\end{document}
